\input{../tex-inputs/latex-leseansicht-vorspann}

\section[Arthur Schnitzler an Gustav Schwarzkopf, {[}4. oder 5.? 11. 1905{]}]{L04397 Arthur Schnitzler an Gustav Schwarzkopf, {[}4. oder 5.? 11. 1905{]}}
\nopagebreak\mylabel{L04397v}
\rehead{ }\normalsize\beginnumbering\briefempfaengerindex{Schwarzkopf, Gustav@\textsc{Schwarzkopf, Gustav}!zzzSchnitzler, Arthur@\emph{von Arthur Schnitzler}!1905-11-051@{{[}4. oder 5.? 11. 1905{]}}|(be}
\toendnotes[C]{\smallbreak\pagebreak[2]}
\correspDesc{Versand  durch Arthur Schnitzler im Zeitraum [4. oder 5.? 11. 1905] in Prag
\newline{}Übermittlung  am 5. [11. 1905?] in Prag
\newline{}Zustellung  am 6. [11. 1905?] in Wien 1
\newline{}Erhalt  durch Gustav Schwarzkopf am 6. [11. 1905?] in Tiefer Graben 23}\toendnotes[C]{\smallbreak}
\Standort{DLA, A:Schnitzler, HS.1985.1.1897.}
\physDesc{Bildpostkarte, 85 Zeichen
\newline{}Handschrift: Bleistift, deutsche Kurrent
\newline{}Versand: 1) Stempel: »\nobreak{}\oindex{Prag@\textbf{Prag}|pwk}Prag 1 Prag, 5. \textcolor{gray}{11}. {[}05{]}\nobreak{}«.   2) Stempel: »\nobreak{}\oindex{I., Innere Stadt@\textbf{I., Innere Stadt}|pwk}\textcolor{gray}{Wien 1}, 6. {[}11. 1905{]}, VIII, Bestellt\nobreak{}«. }\toendnotes[C]{\smallbreak}\pstart{}{\pb}Herrn \textsc{Gustav Schwarzkopf}\pend{}\pstart{}\textsc{Wien I}\oindex{I., Innere Stadt@\textbf{I., Innere Stadt}|pw}\pend{}\pstart{}\textsc{Tiefer Graben 23}\oindex{Wien@\textbf{Wien}!I., Innere Stadt@\textbf{I., Innere Stadt}!Tiefer Graben 23@\textbf{Tiefer Graben 23}|pw}.\pend{}{\bigskip}
\pstart
           \centering{}{\pb}\textcolor{gray}{\textbf{Gruss aus \label{K_L04397-1v}\edtext{Prag\oindex{Prag@\textbf{Prag}|pw}}{\lemma{\textnormal{\emph{Prag}}}\Cendnote{\textnormal{Die Briefmarke hat einen Wert von fünf Heller. Zusammen mit der Tagangabe des
               Poststempels (»5.«) kann damit das Datum der Postkarte auf den 4. oder
                     5. 11. 1905 bestimmt werden, als sich Schnitzler
                     für eine Lesung\eventindex{Deutsches Haus [Prag]@\textbf{Deutsches Haus [Prag]}!Lesung von Der blinde Geronimo, Excentric, 5.11.1905@Lesung von Der blinde Geronimo, Excentric, 5.11.1905|pwkv} (Lesung von Der blinde Geronimo, Excentric, 5.11.1905. In: A. S.: \emph{Kulturveranstaltungen}, \url{https://schnitzler-kultur.acdh.oeaw.ac.at/pmb183265.html})
                     in Prag\oindex{Prag@\textbf{Prag}|pwk} aufhielt.}}}\label{K_L04397-1}!}}\pend
           
\pstart
           \centering{}\textcolor{gray}{\textbf{Das neue deutsche Theater\oindex{Neues Deutsches Theater@\textbf{Neues Deutsches Theater}|pw}. – 86.}}\pend
           \vspace{1em}
\pstart
           \noindent{}{\pb}Herzlichen Gruſs Ihnen und Dr. Max\pwindex{Schwarzkopf, Max 12.\,6.\,1857 Wien – 14.\,4.\,1928 ebd.@\textsc{Schwarzkopf, Max} (12.\,6.\,1857 Wien – 14.\,4.\,1928 ebd.)|pw}.\pend
           
\pstart
           Ihr{\\[\baselineskip]}\spacefill\mbox{A.}\pend
           \leftskip=0em{}\selectlanguage{ngerman}\endnumbering\briefempfaengerindex{Schwarzkopf, Gustav@\textsc{Schwarzkopf, Gustav}!zzzSchnitzler, Arthur@\emph{von Arthur Schnitzler}!1905-11-041@{{[}4. oder 5.? 11. 1905{]}}|)be}\mylabel{L04397h}
\begin{anhang}
\end{anhang}\newcommand{\dateiname}{L04397}\newcommand{\titel}{Arthur Schnitzler an Gustav Schwarzkopf, [4. oder 5.? 11. 1905]}\newcommand{\editorInnen}{Herausgegeben von Jahnke, SelmaMüller, Martin Anton}\input{../tex-inputs/latex-leseansicht-abspann}
